\documentclass[DIN, pagenumber=false, fontsize=11pt, parskip=half]{scrartcl}

\usepackage[utf8]{inputenc}
\usepackage[T1]{fontenc}
\usepackage{booktabs} % for \midrule
\usepackage{template}

\titlehead{\centering\small
    Worked solutions to selected exercises from\\
    Richard Hammack, \emph{Book of Proof, edition 2.2}\\
\rule{0.4\linewidth}{0.4pt}}

\title{}
\author{} % Oscar A. Carrera
\date{}

\begin{document}
\maketitle

\setcounter{section}{2}

\section{Counting}

\subsection{Counting Lists}

\begin{problems}
    \problem Airports are identified with 3-letter codes. For example, the Richmond, Virginia airport has the code \textit{RIC}, and Portland, Oregon has \textit{PDX}. How many different 3-letter codes are possible? 
    \answer There are $26 \cdot 26 \cdot 26 = 17576$ possible 3-letter codes that are different. 

    \problem Five cards are dealt off of a standard 52-card deck and lined up in a row. How many such line-ups are there in which all 5 cards are of the same color (i.e., all black or all red)?
    \answer From a 52-card deck, there are 26 of each color, from each half we have $26 \cdot 25 \cdot 24 \cdot 23 \cdot 22 = 7,893,600$ line-ups where all 5 cards are of that same particular color, taking into consideration it can be either of the two colors, then we have $2 \cdot 7,893,600 = 15,787,200$ line-ups where all 5 cards are of the same color.

    \problem Five cards are dealt off a standard 52-card deck and lined up in a row. How many such line-ups are there in which exactly one of the 5 cards is a queen?
    \answer We have 5 positions to place one of the 4 queens, that is $4 \cdot 5 = 20$ different ways, the remaining 4 positions can be any of the remaining 48 cards that are not a queen, therefore there are $4 \cdot 5 \cdot 48 \cdot 47 \cdot 46 \cdot 45 = 93,398,400$ line-ups with such arrangement.

    \problem This problem concerns lists made from the symbols \textit{A,B,C,D,E}.
    \begin{enumerate}[label=\alph*)]
        \item How many such length-5 lists have at least one letter repeated?
        \item How many such length-6 lists have at least one letter repeated?
    \end{enumerate}
    \answer
    \begin{parts}
        \item We have $5^5$ different possible lists and $5!$ which don't repeat, thus we have $5^5 - 5! = 3005$ length-5 lists with at least one letter repeated.
        \item There are more positions than symbols, thus all possible length-6 lists have at least one letter repeated, so there are $5^6 = 15,625$ such lists which have at least one letter repeated.
    \end{parts}

    \problem This problem concerns lists made from the letters \textit{A,\allowbreak B,\allowbreak C,\allowbreak D,\allowbreak E,\allowbreak F,\allowbreak G,\allowbreak H,\allowbreak I,\allowbreak J}.
    \begin{enumerate}[label=\alph*)]
        \item How many length-5 lists can be made from these letters if repetition is not allowed and the list must begin with a vowel?
        \item How many length-5 lists can be made from these letters if repetition is not allowed and the list must begin and end with a vowel?
        \item How many length-5 lists can be made from these letters if repetition is not allowed and the list must contain exactly one A?
    \end{enumerate}
    \answer
    \begin{parts}
        \item There are 3 vowels for the first position, the remaining 4 positions are filled from the other 9 without repetition, so there are $3 \cdot 9 \cdot 8 \cdot 7 \cdot 6 = 9,072$ such lists.
        \item There are 3 choices of vowel for the first position and 2 for the last, and the middle 3 positions are filled from the remaining 8 letters, so there are $3 \cdot 2 \cdot 8 \cdot 7 \cdot 6 = 2,016$ such lists.
        \item The single \textit{A} can go in any of the 5 positions, and the remaining 4 positions are filled from the 9 letters that are not \textit{A}, so there are $5 \cdot 9 \cdot 8 \cdot 7 \cdot 6 = 15,120$ such lists.
    \end{parts}

    \problem Consider the lists of length six made with the symbols \textit{P,R,O,F,S}, where repetition is allowed. (For example, the following is such a list: (P,R,O,O,F,S).) How many such lists can be made if the list must end in an S and the symbol O is used more than once?
    \answer The last position is fixed as \textit{S}, leaving the first 5 positions free, so there are $5^5 = 3125$ lists ending in \textit{S}. From these we subtract the ones where \textit{O} appears fewer than two times: $4^5 = 1024$ lists with no \textit{O}, and $5 \cdot 4^4 = 1280$ lists with exactly one \textit{O} (5 positions for the \textit{O}, the other 4 filled from the 4 remaining symbols). Thus there are $5^5 - 4^5 - 5 \cdot 4^4 = 821$ such lists.

\end{problems}

\subsection{Factorials}

\begin{problems}
    \problem For which values of $n$ does $n!$ have $n$ or fewer digits?
    \answer We assume the problem refers to the length of the list of digits used to represent $n!$ for some $n$ in base 10. If we let $D(n)$ represent the length of the list for such representation, then the problem statement can be written as $D(n) \leq n$.

    From the \textit{Basis Representation Theorem} we know we can represent any positive integer $N$ uniquely in the form $N = \sum_{k=0}^{m}a_k10^k$, $a_m \neq 0, a_k \in \set{0,...,9}$. From this, we see the number of digits of $N$ is exactly $m+1$, and also $10^m \leq N \leq 10^{m+1}-1$.

    If we let $N = n!$, and $d$ be the number of digits of $n!$, then $10^{d-1} \leq n! \leq 10^d - 1 < 10^d$, so $d - 1 \leq \log_{10}(n!) < d$ which is equivalent to $d = \lfloor \log_{10}(n!) \rfloor + 1$. Therefore we have $D(n) = \lfloor \log_{10}(n!) \rfloor + 1$.

    Using \textit{Stirling's formula}, if we let $\delta$ account for some error, then we have, $\log_{10}(n!) = \log_{10}(\sqrt{2\pi n}(\frac{n}{e})^n) + \delta \approx n\log_{10}(\frac{n}{e}) + \frac{1}{2}\log_{10}(2\pi n)$.

    We can find an approximation to the \textit{crossover} point where $\log_{10}(n!)$ grows like $n\log_{10}(\frac{n}{e})$ that is, $n\log_{10}(\frac{n}{e}) = n$, thus,
    \begin{align}
        \nonumber n\log_{10}\left(\frac{n}{e}\right) & = n \\
        \nonumber \log_{10}\left(\frac{n}{e}\right) & = 1\\
        \nonumber \frac{n}{e} & = 10\\ 
        \nonumber n & = 10e \approx 27.18
    \end{align}
    We try some discrete neighbor values:
    \begin{align}
        \nonumber 28! = 304,888,344,611,713,860,501,504,000,000 & \rightarrow 30  \text{ digits}\\ 
        \nonumber 27! = 10,888,869,450,418,352,160,768,000,000 & \rightarrow 29 \text{ digits}\\ 
        \nonumber 26! = 403,291,461,126,605,635,584,000,000 & \rightarrow 27 \text{ digits}\\ 
        \nonumber 25! = 15,511,210,043,330,985,984,000,000 & \rightarrow 26 \text{ digits}\\
        \nonumber 24! = 620,448,401,733,239,439,360,000 & \rightarrow 24 \text{ digits}
    \end{align}
    We see that when $n = 24$, then $n!$ has $n$ or fewer digits, and the trend continues for $1 \leq n \leq 24$.


    \problem Using only pencil and paper, find the value of $\frac{100!}{95!}$.
    \answer $\frac{100!}{95!} = \frac{100\cdot99\cdot98\cdot97\cdot96\cdot95\cdot94\cdot93\cdot...}{95\cdot94\cdot93\cdot...} = 100\cdot99\cdot98\cdot97\cdot96 = 9034502400$

    \problem There are two 0's at the end of $10! = 3,628,800$. Using only pencil and paper, determine how many 0's are at the end of the number $100!$. 
    \answer The question in hand is determining how many times is $100!$ divisible by 10, in other words, if $100! = 2^m \cdot 5^n \cdot \delta$ (where $5 \nmid \delta$ and $2 \nmid \delta$), then each pair of one 2 and one 5 makes a factor of 10:
    \[
        2^m \cdot 5^n = 10^{\min(m,n)} \cdot 2^{m-\min(m,n)} \cdot 5^{n-\min(m,n)}
    \]
    Since $100!$ contains more factors of 2 than factors of 5, then $m > n$, thus:
    \[
        100! = 10^n \cdot 2^{m-n} \cdot \delta
    \]
    So the number of trailing zeros is exactly $n$, and we obtain $n$ by counting the multiples of 5 in $1 \ldots 100$:
    \[
        \frac{100}{5} = 20
    \]
    So there is at least 20 trailing zeroes, but we have to account for factors that contribute more than one factor, which is the case of 25, 50, 75, 100:
    \begin{align}
        \nonumber 25 & = 5^2\\ 
        \nonumber 50 & = 2 \cdot 5^2\\ 
        \nonumber 75 & = 3 \cdot 5^2\\ 
        \nonumber 100 &= 4 \cdot 5^2
    \end{align}
    Hence we need to account for $5^2$'s:
    \[
        \frac{100}{25} = 4
    \]
    Then we have that $100! = 2^m \cdot 5^{20+4} \cdot \delta \implies n = 20 + 4 = 24$, that is, there are 24 trailing zeroes.


    \problem Compute how many 7-digit numbers can be made from the digits 1,2,3,4,5,6,7 if there is no repetition and the odd digits must appear in an unbroken sequence.
    \answer There are $4! \cdot 4! = 576$ such numbers.

    \problem There is another significant function called \textbf{Stirling's formula} that provides an approximation to factorials. It states that $n! \approx \sqrt{2\pi n}(\frac{n}{e})^n$. It is an approximation to $n!$ in the sense that $\frac{n!}{\sqrt{2 \pi n }(\frac{n}{e})^n}$ approaches 1 as $n$ approaches $\infty$. Use Stirling's formula to find approximations to $5!, 10!, 20!$ and $50!$.
    \answer Writing $S(n) = \sqrt{2\pi n}\left(\frac{n}{e}\right)^n$ for Stirling's approximation, we obtain:
    \begin{align}
        \nonumber S(5)  & \approx 118.02\\
        \nonumber S(10) & \approx 3{,}598{,}695.62\\
        \nonumber S(20) & \approx 2.4228 \times 10^{18}\\
        \nonumber S(50) & \approx 3.0363 \times 10^{64}
    \end{align}
\end{problems}

\subsection{Counting Subsets}

\begin{problems}
    \problem Suppose $A$ is a set for which $|A| = 100$. How many subsets of $A$ have 5 elements? How many subsets have 10 elements? How many have 99 elements?
    \answer
    \begin{align}
        \nonumber {100 \choose 5} & = \frac{100!}{5!(100-5)!} & = 75,287,520\\ 
        \nonumber {100 \choose 10} & = \frac{100!}{10!(100-10)!} & = 17,310,309,456,440\\ 
        \nonumber {100 \choose 99} & = \frac{100!}{99!(100-99)!} & = 100
    \end{align}


    \problem Suppose a set $B$ has the property that $|\set{X : X \in \mathcal{P}(B),|X| = 6}| = 28$. Find $|B|$.
    \answer From the problem statement we know that ${\abs{B} \choose 6} = 28$, since 28 is small, we check neighbor values of 6:
    \begin{align}
        \nonumber {6 \choose 6} & = 1\\ 
        \nonumber {7 \choose 6} & = 7\\ 
        \nonumber {8 \choose 6} & = 28
    \end{align}
    So we see that $\abs{B} = 8$.

    \problem $|\set{X \in \mathcal{P}(\set{0,1,2,3,4,5,6,7,8,9}) : |X| = 4}| = $
    \answer ${10 \choose 4} = 210$

    \problem This problem concerns lists made from the symbols $A,\allowbreak B,\allowbreak C,\allowbreak D,\allowbreak E,\allowbreak F,\allowbreak G,\allowbreak H,\allowbreak I,\allowbreak J$.
    \begin{enumerate}[label=\alph*)]
        \item How many length-5 lists can be made if repetition is not allowed and the list is in alphabetical order? (Example: \textit{BDEFI} or \textit{ABCGH}, but not \textit{BADGH}.)
        \item How many length-5 lists can be made if repetition is not allowed and the list is \textbf{not} in alphabetical order?
    \end{enumerate}
    \answer
    \begin{parts}
    \item There are ${10 \choose 5} = 252$ ways to choose 5 letters from 10.
    \item Since repetition is not allowed, there are $\frac{10!}{5!} = 30,240$ such lists, from the previous question we know there are 252 lists in alphabetical order, therefore, there are $30,240 - 252 = 29,988$ such lists.
    \end{parts}

    \problem A department consists of 5 men and 7 women. From this department you select a committee with 3 men and 2 women. In how many ways can you do this?
    \answer There are ${5 \choose 3} \cdot {7 \choose 2} = 210$ different ways to create the committee.

    \problem Twenty-one people are to be divided into two teams, the Red Team and the Blue Team. There will be 10 people on Red Team and 11 people on Blue Team. In how many ways can this be done?
    \answer This can be done ${21 \choose 10} = 352,716$ different ways.

    \problem Suppose $n,k \in \mathbb{Z}$, and $0 \leq k \leq n$. Use Definition 3.2 alone (without using Fact 3.3) to show that ${n \choose k}  = {n \choose n-k}$.
    \answer Let $S$ be a set with $|S| = n$. By Definition 3.2, ${n \choose k}$ is the number of $k$-element subsets of $S$, and ${n \choose n-k}$ is the number of $(n-k)$-element subsets of $S$.

    Consider pairing each subset $A \subseteq S$ with its complement $S - A$. If $|A| = k$, then $|S - A| = n - k$, so this pairs every $k$-element subset with an $(n-k)$-element subset. Furthermore, taking the complement twice returns the original set, since $S - (S - A) = A$; hence each $(n-k)$-element subset is paired back with exactly one $k$-element subset, with none left out or counted twice.

    So the $k$-element subsets and the $(n-k)$-element subsets are in exact one-to-one correspondence, and therefore equal in number. By Definition 3.2, $\binom{n}{k} = \binom{n}{n-k}$.
\end{problems}

\subsection{Pascal's Triangle and the Binomial Theorem}

\begin{problems}
    \problem Use the binomial theorem to find the coefficient of $x^8y^5$ in $(x+y)^{13}$.
    \answer $(x + y)^{13} = ... + {13 \choose 5}x^8y^5 + ... = ... + 1287x^8y^5 + ...$

    \problem Use the binomial theorem to find the coefficient of $x^6y^3$ in $(3x-2y)^9$.
    \answer $(3x - 2y)^9 = (3x + (-2y))^9 = ... + {9 \choose 3}(3x)^{6}(-2y)^3 + ... = ... - 489888x^6y^3 + ...$

    \problem Use Definition 3.2 and Fact 1.3 to show $\sum_{k=0}^{n} {n \choose k} = 2^n$.
    \answer By Definition 3.2, ${n \choose k}$ is the number of subsets that can be made by choosing $k$ elements from a set with $n$ elements, and from Fact 1.3 we know a finite set with $n$ elements has $2^n$ subsets.

    $\sum_{k=0}^{n}{n \choose k}$ is the sum of all possible subsets, and since each of the $n$ elements can be either included or not included in a subset, there are $2 \cdot 2 \cdot ... \cdot 2 = 2^n$ possible subsets.

    Therefore, $\sum_{k=0}^{n}{n \choose k} = 2^n$

    \problem Use Fact 3.3 to derive Equation 3.2.
    \answer Fact 3.3 states that ${n \choose k} = \frac{n!}{k!(n-k)!}$ for $n,k \in \mathbb{Z}$ and $0 \leq k \leq n$. Equation 3.2 states that ${n + 1 \choose k} = {n \choose k - 1} + {n \choose k}$.
    \begin{align}
        \nonumber {n \choose k - 1} + {n \choose k} & = \frac{n!}{(k-1)!(n-k+1)!} + \frac{n!}{k!(n-k)!}\\ 
        \nonumber & = \frac{n!k}{k!(n-k+1)!} + \frac{n!(n-k+1)}{k!(n-k+1)!}\\ 
        \nonumber & = \frac{n!k + n!(n-k+1)}{k!(n-k+1)!} = \frac{n!(n+1)}{k!(n-k+1)!}\\
        \nonumber & = \frac{(n+1)!}{k!(n + 1 - k)!}\\ 
        \nonumber & = {n + 1 \choose k}
    \end{align}

    \problem Show that the formula $k{n \choose k} = n{n-1 \choose k-1}$ is true for all integers $n,k$ with $0 \leq k \leq n$.
    \answer
    \begin{align}
        \nonumber n {n-1 \choose k-1} & = n\left(\frac{(n-1)!}{(k-1)!([n-1]-[k-1])!}\right)\\ 
        \nonumber & = n\left(\frac{(n-1)!}{(k-1)!(n-k)!}\right)\\ 
        \nonumber & = n\left(\frac{k(n-1)!}{k!(n-k)!}\right)\\ 
        \nonumber & = \frac{kn(n-1)!}{k!(n-k)!}\\ 
        \nonumber & = k\left(\frac{n(n-1)!}{k!(n-k)!}\right)\\ 
        \nonumber & = k\left(\frac{n!}{k!(n-k)!}\right)\\
        \nonumber & = k{n \choose k}
    \end{align}

    \problem Show that ${n \choose k} {k \choose m} = {n \choose m} {n-m \choose k-m}$.
    \answer 
    \begin{align}
        \nonumber {n \choose k}{k \choose m} & = \left(\frac{n!}{k!(n-k)!}\right) \left(\frac{k!}{m!(k-m)!}\right)\\ 
        \nonumber & = \frac{n!}{m!(k-m)!(n-k)!}\\ 
        \nonumber {n \choose m}{n-m \choose k-m} & = \left(\frac{n!}{m!(n-m)!}\right) \left(\frac{(n-m)!}{(k-m)!([n-m]-[k-m])!}\right)\\ 
        \nonumber & = \frac{n!}{m!(k-m)!(n-k)!}
    \end{align}


    \problem The first five rows of Pascal's triangle appear in the digits of powers of 11: $11^0 = 1$, $11^1 = 11$, $11^2 = 121$, $11^3 = 1331$, $11^4 = 14641$. Why is this so? Why does the pattern not continue with $11^5$?
    \answer This comes from the Binomial Theorem:

    $(10 + 1)^n = \sum_{k=0}^{n}{n \choose k}10^{n-k}$

    And as long every binomial coefficient is a single digit, the decimal representation of $11^n$ is simply the row of Pascal's triangle.

    This fails at $11^5$ because the 5th row of Pascal's triangle contains coefficients with double digits.
\end{problems}

\end{document}
